\documentclass[12pt]{article}
\usepackage[a4paper,margin=2.2cm]{geometry}
\usepackage{amsmath,amssymb,amsthm,mathtools}
\usepackage{hyperref}
\usepackage{cleveref}

% 中文（xelatex）
\usepackage{fontspec}
\usepackage{xeCJK}

\usepackage{fontspec}
\usepackage{xeCJK}
\setCJKmainfont{Noto Serif CJK SC}

\title{QR}
\date{}
\begin{document}
\maketitle
\begin{flushleft}

    1: QR 前置知识: \\
    由 Least-Squares 我们得到关键等式: \\
    \begin{align}
        A^{T}A\hat{c} = A^{T}u \\
        A^{T}p = A^{T}A\hat{c}
    \end{align}

    在 $\mathbb{R}^{3}$ 空间中 存在 W 子空间 $W = span(a_1,a_2)\; (a_1, a_2) \in \mathbb{R}^{3 \times 1}, Col(A) = W = [a_1, a_2]$ \\

    我们这里回顾下 $A$ 和 $A{T}$: \\
    1: $A, 和 A^{T}$ 都是矩阵, 是转置关系, $A \in \mathbb{R}^{3 \times 2}, A^{T} \in \mathbb{R}^{2 \times 3}$ \\
    2: $Col(A) \subseteq \mathbb{R}^{3}, Col(A^{T}) \subseteq \mathbb{R}^{2}$ $(Rank \leq min(rows_n,cols_n) \land \text{列之间线性无关})$ \\
    3: $Col(A)$ 的基底组是: $[a_1, a_2]$
    4: $Col(A^{T})$ 的基底组是什么?  \\
    已知 $Col(A)$ 的基底组是 $\begin{bmatrix}
            a_{11} & a_{21} \\
            a_{12} & a_{22} \\
            a_{13} & a_{23} \\
        \end{bmatrix}$ \\
    $Col(A^{T})$ = $\begin{bmatrix}
            a_{11} & a_{12} & a_{13} \\
            a_{21} & a_{22} & a_{23} \\
        \end{bmatrix}$
    已知道 $Rank(Col(A^{T})) = 2$ 那么 $Col(A^{T})$ 的列必然存在一组线性相关
    我们假设$col1$与$col2$线性无关(因为 rank = 2, 所以能保证能找到两列独立列),
    那么 $Col(A^{T})$的基底组是: $\begin{bmatrix}
            a_{11} & a_{12} \\
            a_{21} & a_{22} \\
        \end{bmatrix}$

    我们再分析等式: $A^{T}p = A^{T}A\hat{c}$ \\
    已知 $p = proj_{w}(u) \in W$, $A\hat{c} \in W$ \\
    $A^{T}A\hat{c}$ 很容易理解: 把 $A\hat{c}$ 点映射到 $Col(A^{T})$ 空间 \\
    那么映射到  $Col(A^{T})$ 空间之后 我们会得到一个 $Col(A^{T})$的向量 z, 也就是$z = A^{T}A\hat{c} = A^{T}p$ \\
    我们展开 z 的等式: \\
    \begin{align}
        z = A^{T}p = \begin{bmatrix}
                         a_1^{T} \\
                         a_2^{T} \\
                     \end{bmatrix}p \\
        =\begin{bmatrix}
             <a_1, p> \\
             <a_2, p> \\
         \end{bmatrix}
    \end{align}
    这里有一个最重要的点 z 到底是一个分量向量(也就是$Col(A^{T})$的基向量的单位数) 但还是单纯的标量向量? 它是单纯的标量向量(占比比例), 它不是为了确定在$Col(A^{T})$的位置。 \\

    我们这里再次问自己一个问题: 最小二乘(Least-Squares) 到底是要解决什么问题 ? \\
    它解决的是无解的情况下的近似解(降维近似解)的问题: \\
    超定线性系统:
    \begin{align}
        Ax = b
    \end{align}
    A: $m \times n$ (m 个方程, n 个未知量, m > n) \\
    x: n 维未知数向量 \\
    b: m 维常数项向量 \\
    对于超定线性系统 当约束条件(方程) > 自由度（未知数）时  \\
    对于 $\vec{x}$ 它永远都固定在 $\mathbb{R}^{n}$  \\
    对于 $\vec{b}$ 和 $Col(A)$ 它属于 $\mathbb{R}^{m}$ \\
    这里超定线性系统 会出现两种结果情况: \\
    1: 无解: \\
    $\vec{b} \notin Col(A)$,不存在一个点 x($x \in \mathbb{R}^{2}$) 使等式成立 \\
    2: 只有当多出来的方程完全是其他方程的线性组合（即多余的方程是冗余的），且常数项 b 也刚好吻合时，才存在解(这不在本次的讨论范围)

    Least-Squares 就是为了解决这个问题的工具, 它使用了几何的正交投影给出一个最佳近似的低维解(m 映射 n)
    分析核心公式： \\
    \begin{align}
        A^{T}A\hat{c} = A^{T}u
    \end{align}

    $u$ 是一个高维观测数据 $A$ 是你希望的低维环境 $A{T}$ 是对应的高维解映射到 $A$的矩阵运算(高维向量转化为低维的内积向量), 其实$A\hat{c}= u$是无解的 \\
    但 基于 $\hat{c}$(Least-Squares) 的 $A^{T}A\hat{c}$ 是有对应的解的($\hat{c}$ 就是低维解 它是存在的, $A^{T}$ 是一个观测空间 证明等式是存在的)
    它是 $A{T}u$的最佳近似的低维解(在 $Col(A^{T})$ 观察子空间) \\

    这里我们要再次的申明 $A^{T}$ 的声明与作用(这很关键): \\
    1:$A^{T}$ 与 $u$ \\
    存在下列代数式子: \\
    \begin{align}
        u = p + e                                                       \\
        e \perp Col(A) \land p \in Col(A)                               \\
        \sum_{i = 1}^{n} <a_{i},e> = A^{T}e= 0                          \\
        A^{T}u = A^{T}(p + e) = A^{T}p + A^{T}e  =  A^{T}p + 0 = A^{T}p \\
    \end{align}
    我们发现 e 残差向量被去掉了 得到一个 $z  = A^{T}p$  z 是内积标量向量 是一个内积测量向量 在等式的一边 检查 $\hat{c}$ 是不是 这个等式的解
    也就是 $A^{T}$ 是一个测量空间也是一个过滤空间(去噪)

    2: QR- Measurement Cast to Coordinates
    分析 $z = A^{T}A\hat{c}$ \\
    已知道 $A = \begin{bmatrix}a_1 & a_2\end{bmatrix} ((a_1,a_2) \in \mathbb{R}^{3 \times 1})$ \\
    $if: <a_1,a_2> = 0$($a_1, a_2$ 正交) \\
    \begin{align}
        A^{T} = \begin{bmatrix}
                    a_1^{T} \\
                    a_2^{T}
                \end{bmatrix}                \\
        A^{T}A = \begin{bmatrix}
                     a_1^{T} \\
                     a_2^{T}
                 \end{bmatrix} \begin{bmatrix}
                                   a_1 & a_2
                               \end{bmatrix} \\
        = \begin{bmatrix}
              a_1^{T}a_1 & a_1^{T}a_2 \\
              a_2^{T}a_1 & a_2^{T}a_2
          \end{bmatrix}             \\
        = \begin{bmatrix}
              ||a_1||^{2} & 0           \\
              0           & ||a_2||^{2}
          \end{bmatrix}
    \end{align}
    我们对 $a_1, a_2$ 单位化
    $ \begin{bmatrix}
            ||a_1||^{2} & 0           \\
            0           & ||a_2||^{2}
        \end{bmatrix}  = \begin{bmatrix}
            1 & 0 \\
            0 & 1
        \end{bmatrix} $ \\
    那么 $z = A^{T}A\hat{c} = I\hat{c} = \hat{c}$ \\
    else: \\
    能找到一组与 A 张成相同子空间、但标准正交的基
    $Q = \begin{bmatrix}q_1 & q_2\end{bmatrix} ((q_1,q_2) \in \mathbb{R}^{3 \times 1})$  \\
    $Q^{T}Q = I$ \\
    替换: \\
    $z = Q^{T}Qy = Iy = y$ \\
    $Qy = p$ \\
    我们就可以说 z 就是 p 在正交坐标系 Q 的坐标

    3: 引入 QR \\
    $A \to QR$ \\
    $A\hat{c} = QR\hat{c}$ \\
    $Col(A) = Col(Q) = W$ \\
    $c \to Rc \to y \to Qy \to p$ \\
    \begin{align}
        A^{T}A\hat{c} = A^{T}u                               \\
        (QR)^{T}(QR)\hat{c} = (QR)^Tu                        \\
        R^{T}Q^{T}QR\hat{c} = R^{T}Q^{T}u                    \\
        R^{T}IR\hat{c} = R^{T}Q^{T}u                         \\
        (R^{T})^{-1}R^{T}IR\hat{c} = (R^{T})^{-1}R^{T}Q^{T}u \\
        IR\hat{c} = Q^{T}u                                   \\
        R\hat{c} = Q^{T}u
    \end{align}
    代数整理之后我们得到了 $R\hat{c} = z$ 这个 z 是 $Col(Q)$ 的标量向量 \\
    也就是说 $Q^{T}u$ u 通过 $Q^{T}$ 映射出来 内积标量向量 就是 $Col(Q)$ 中的投影坐标(把复杂的投影运算转化为直接做内积（提取坐标）) \\

    4: 计算 \\
    $a_{1} = [1,1,0]^{T}, a_{2} = [1,1,1]^{T}, A = \begin{bmatrix}
            1 & 1 \\
            1 & 1 \\
            0 & 1
        \end{bmatrix}$
    通过 Gram-Schmidt 构造 Q \\
    step1: $w1 = a1$ \\
    归一:  $q1 = \frac{w1}{||w1||}$ \\
    \begin{align}
        q1 = \frac{[1,1,0]^{T}}{\sqrt{1^{2} + 1^{2}}}
        = \frac{[1,1,0]^{T}}{\sqrt{2}}
        = \frac{1}{\sqrt{2}}[1,1,0]^{T}
        = [\frac{1}{\sqrt{2}}, \frac{1}{\sqrt{2}}, 0]^{T}
    \end{align}
    step2: q2  \\
    $w2 = a_{2} - proj_{q1}(a_{2})$ \\
    已知 q1 是单位向量 \\
    \begin{align}
        proj_{q1}(a_{2}) = <q1,a_{2}>q1                                                                         \\                                                                         \\
        = (\frac{1}{\sqrt{2}} * 1 + \frac{1}{\sqrt{2}} * 1 + 0) [\frac{1}{\sqrt{2}}, \frac{1}{\sqrt{2}}, 0]^{T} \\
        = \sqrt{2} [\frac{1}{\sqrt{2}}, \frac{1}{\sqrt{2}}, 0]^{T}                                              \\
        = [1, 1, 0]^{T}                                                                                         \\
    \end{align}
    $w2 =  [1,1,1]^{T} - [1, 1, 0]^{T} = [0, 0, 1]^{T}$ \\
    归一: $q2 = \frac{w2}{||w2||}$ \\
    \begin{align}
        q2 = \frac{[0, 0, 1]^{T}}{1} \\
        q2 = [0, 0, 1]^{T}
    \end{align}
    step3 : Q \\
    \begin{align}
        <q1,q2> = 0                                                                     \\
        = \frac{1}{\sqrt{2}} * 0  + \frac{1}{\sqrt{2}} * 0 +  0 * 1 = 0                 \\                                                         \\
        Q = \begin{bmatrix}
                \frac{1}{\sqrt{2}} & 0 \\
                \frac{1}{\sqrt{2}} & 0 \\
                0                  & 1
            \end{bmatrix}                                        \\
        Q^{T}Q =\begin{bmatrix}
                    \frac{1}{\sqrt{2}} & \frac{1}{\sqrt{2}} & 0 \\
                    0                  & 0                  & 1
                \end{bmatrix}\begin{bmatrix}
                                 \frac{1}{\sqrt{2}} & 0 \\
                                 \frac{1}{\sqrt{2}} & 0 \\
                                 0                  & 1
                             \end{bmatrix} \\
        = \begin{bmatrix}
              \frac{1}{\sqrt{2}} * \frac{1}{\sqrt{2}} + \frac{1}{\sqrt{2}}* \frac{1}{\sqrt{2}} + 0 & 0 \\
              0                                                                                    & 1 \\
          \end{bmatrix}
        = \begin{bmatrix}
              1 & 0 \\
              0 & 1
          \end{bmatrix}
    \end{align}
    延展: 关于$QQ^{T}$ \\
    已知 $Col(Q) = W$, 我们设代数等式: $ k = QQ^{T}u$ \\
    $Q^{T}u$ 表示 u 会得到一个内积比例向量, 然后在进行 Q 矩阵运算 \\
    所以我们能得到 k 就是 u 在 $Col(Q)$ 的投影($k \in \mathbb{R}^{3 \times 1}  \land  k \in Col(Q)$)\\

    step3: 构建 R \\
    \begin{align}
        A = QR                           \\
        Q^{T}A  = Q^{T}QR                \\
        Q^{T}A = R                       \\
        R  = \begin{bmatrix}
                 q_1^{T} \\
                 q_2^{T}
             \end{bmatrix}\begin{bmatrix}
                              a_1, a_2
                          \end{bmatrix} \\
        = \begin{bmatrix}
              <q_1^{T},a_1> & <q_1^{T},a_2> \\
              <q_2^{T},a_1> & <q_2^{T},a_2>
          \end{bmatrix}  \\
        q_2 \perp a_1                    \\
        = \begin{bmatrix}
              <q_1^{T},a_1> & <q_1^{T},a_2> \\
              0             & <q_2^{T},a_2>
          \end{bmatrix}
        = \begin{bmatrix}
              (\frac{1}{\sqrt{2}} + \frac{1}{\sqrt{2}}) & (\frac{1}{\sqrt{2}} + \frac{1}{\sqrt{2}}) \\
              0                                         & 1
          \end{bmatrix}
        = \begin{bmatrix}
              \sqrt{2} & \sqrt{2} \\
              0        & 1        \\
          \end{bmatrix}
    \end{align}
    验证: $QR = A$ \\
    \begin{align}
        \begin{bmatrix}
            \frac{1}{\sqrt{2}} & 0 \\
            \frac{1}{\sqrt{2}} & 0 \\
            0                  & 1
        \end{bmatrix} \begin{bmatrix}
                          \sqrt{2} & \sqrt{2} \\
                          0        & 1        \\
                      \end{bmatrix}                                                                \\
        = \begin{bmatrix}
              \frac{1}{\sqrt{2}} * \sqrt{2} & \frac{1}{\sqrt{2}} * \sqrt{2} + 0 * 1 \\
              \frac{1}{\sqrt{2}} * \sqrt{2} & \frac{1}{\sqrt{2}} * \sqrt{2} + 0 * 1 \\
              0                             & 1                                     \\
          \end{bmatrix} \\
        = \begin{bmatrix}
              1 & 1 \\
              1 & 1 \\
              0 & 1 \\
          \end{bmatrix}
    \end{align}

    5: 我们来思考 QR 到底解决了什么问题(简化了什么问题) \\
    已证下列关键代数式: \\
    \begin{align}
        A^{T}A\hat{c} = A^{T}u                             \\
        (A^{T}A)^{-1}(A^{T}A)\hat{c} = (A^{T}A)^{-1}A^{T}u \\
        \hat{c} = (A^{T}A)^{-1}A^{T}u
    \end{align}
    也就是说我们要算出 Least-Squares 我们必须要进行复杂的逆矩阵运算 \\
    引入 QR 之后 \\
    \begin{align}
        R\hat{c} = Q^{T}u             \\
        R^{-1}R\hat{c} = R^{-1}Q^{T}u \\
        \hat{c} = R^{-1}Q^{T}u
    \end{align}



\end{flushleft}

\end{document}